\chapter{\texorpdfstring{$p$}{p}-adic CY Langlands for the Fermat quartic}
\label{ch:padic-langlands-cy}

The chiral image $\Phi_2(X) = \SpCh_{\Sigma_1, C}(\PhiFA_2(D^b\Coh(X)))$ of a compact CY$_2$ surface $X$ is the Mukai Heisenberg $H_{\mathrm{Muk}}$ of rank $24 = b_2(X) + 2 = \kappa_{\mathrm{fiber}}(X)$ (CY-A$_2$; two-stage factorisation Theorem~\ref{thm:phi-two-stage-factorisation} at $d = 2$, $\Etwo$-native by Proposition~\ref{prop:native-en-level}). The $p$-adic deformations land on outputs of $\PhiFA_2$: the Frobenius action on $H^\ast(X, \Q_\ell)$ deforms the specialisation datum $(\Sigma_1, C)$ over $\Z_p$, with $\PhiFA_2(D^b\Coh(X))$ the stable geometric input. When $X$ is defined over $\Q$ and carries complex multiplication on its transcendental lattice, the $p$-adic shadow zeta of $\Phi_2(X)$ matches the Hasse--Weil $L$-function of $X$ and, via the Kuga--Satake construction, a classical weight-$3$ CM newform. The first proved instance is the Fermat quartic $x_0^4 + x_1^4 + x_2^4 + x_3^4 = 0$, a singular K3 surface over $\Q$ of Picard rank $\rho_{\mathrm{NS}} = 20$ whose transcendental lattice carries complex multiplication by $\Q(i)$. The proof assembles CY-A$_2$ (chiral side), Livne (1995) modularity (Galois-representation side), and Kuga--Satake (1967) (Hodge-theoretic bridge); the genus-one shadow $\Theta_1(\Phi_2(X))$ localises the motivic content in the bar complex of $H_{\mathrm{Muk}}$.

\begin{remark}[Stage discipline: Frobenius deforms stage~\texorpdfstring{$2$}{2}; \texorpdfstring{$\PhiFA_2$}{PhiFA-2} is the stable stage~\texorpdfstring{$1$}{1} input]
\label{rem:padic-stage-discipline}
The arithmetic inputs of this chapter sort cleanly across the two stages
of $\Phi_2 = \SpCh_{\Sigma_1, C} \circ \PhiFA_2$.
\begin{enumerate}[label=\textup{(\roman*)}]
\item \textbf{Stage~$1$, stable geometric input.} The canonical object
$\PhiFA_2(D^b\Coh(X)) \in \mathrm{E}_2\text{-}\mathrm{hFA}(X)$ is the
\emph{rigid} piece of the $p$-adic correspondence: its existence does not
deform under Frobenius, and its rank-$24$ Mukai-Heisenberg shape is
independent of $p$ and of any choice of $(\Sigma_1, C)$. The
tier-(i) invariants -- Mukai lattice rank $24$, Hodge numbers
$h^{0,q}(X) = (1, 0, 1)$ -- and the tier-(ii) invariant
$\kappa_{\mathrm{fiber}}(X) = 24$ are all stage-$1$ data.
\item \textbf{Stage~$2$, Frobenius-deformed specialisation.} The Frobenius
action on $H^\ast(X, \Q_\ell)$ deforms the stage-$2$ specialisation
datum $(\Sigma_1, C)$ over $\Z_p$: different primes of good reduction
produce distinct $(\Sigma_1, C)_{\F_p}$ data, distinct chiral shadows
$\Phi_2(X)_{\F_p}$, and distinct Frobenius eigenvalues on $H^2(X)$. The
$p$-adic shadow zeta $\zeta^{(p)}_{\Phi(X)}$ and the Hasse--Weil local
factor $L_p(H^2(X), s)$ are both stage-$2$ quantities indexed by
$p \in \{3, 5, 7, \ldots\}$; the bridge identity $a_n(p) =
1 + \Tr(\Frob_p^n \mid H^2(X)) + p^{2n}$
(Lemma~\ref{lem:phi-k3-bar-shadow-hasse-weil-n1},
eq.~\eqref{eq:shadow-l-bridge}) is a stage-$2$ identification of two
readings of the $\Frob_p$-spectrum.
\end{enumerate}
Livne modularity
(Proposition~\ref{prop:phi-k3-padic-langlands-fermat}(ii)) is a statement
about the stage-$2$ Galois representation on $T(X)$; Kuga--Satake
recovery is the Hodge-theoretic bridge between stage~$1$
($\HH^\bullet(\cC)$-input) and stage~$2$ (Frobenius-deformed
specialisation). This structural separation is what allows the
$p$-adic CY Langlands identity to be stated uniformly in $p$: the
$\PhiFA_2$ stage is stable, and all prime-dependent content sits on the
$\SpCh$ stage.
\end{remark}


\section{The Fermat quartic and the Mukai Heisenberg}
\label{sec:padic-fermat-phi}

The Fermat quartic $X \subset \P^3_\Q$ is a K3 surface over $\Q$ with geometric Picard rank $\rho_{\mathrm{NS}} = 20$ and transcendental lattice $T(X)$ of rank $2$. The Shioda--Inose structure identifies $T(X)$ with the period lattice of an elliptic curve $E_i$ of $j$-invariant $1728$, whose endomorphism ring is $\Z[i]$; accordingly $T(X) \otimes \Q_\ell$ carries complex multiplication by $\Q(i)$. The bad primes of $X / \Q$ are concentrated at $p = 2$, where $X$ has additive reduction; the minimal newform level for the associated Galois representation is $16 = 2^4$.

The CY-to-chiral image of $X$ at $d = 2$ is the rank-$24$ Mukai Heisenberg:
\[
 \Phi(X) \;=\; H_{\mathrm{Muk}}(X),
 \qquad \kappa_{\mathrm{ch}}^{\mathrm{Hodge}}(\Phi(X)) \;=\; \chi(\cO_X) \;=\; 2,
 \qquad \kappa_{\mathrm{fiber}}(\Phi(X)) \;=\; 24.
\]
This identification is the $d = 2$ content of CY-A$_2$.


\section{The lemma on \texorpdfstring{$p = 1$}{p = 1} Dirichlet coefficients}
\label{sec:padic-fermat-lemma}

The higher Dirichlet coefficients require the full Frobenius spectrum.
At $n = 1$ no approximation is needed; the clean statement is the
following lemma, which is load-bearing for the proposition.

\begin{lemma}[Fermat-quartic first Dirichlet coefficient]
\label{lem:phi-k3-bar-shadow-hasse-weil-n1}
Let $X \subset \P^3_\Q$ be the Fermat quartic and $\Phi(X) = H_{\mathrm{Muk}}$ the CY-to-chiral image at $d = 2$. For every rational prime $p \neq 2$ of good reduction, the $p$-adic shadow-zeta Dirichlet coefficient at $n = 1$ satisfies
\[
 a_1(p) \;=\; 1 + \Tr(\Frob_p \mid H^2(X, \Q_\ell)) + p^2.
\]
This equality is the specialisation $\mathbb{L} = p^2$ of the CY-A$_2$ local-factor identity at weight $n = 1$ and uses only $H^0(X) \cong \Q_\ell$, $H^4(X) \cong \Q_\ell(-2)$, and $H^1(X) = H^3(X) = 0$.
\end{lemma}

\begin{proof}
By CY-A$_2$ the motivic bar dimensions of $\Phi(X)$ at charge $n$,
refined by the Frobenius action on $H^\ast(X, \Q_\ell)$ and evaluated at
$\mathbb{L} = p^2$, produce the shadow-zeta Dirichlet coefficients of
$\Phi(X)$. At charge $n = 1$, $a_1(p) = \Tr(\Frob_p \mid H^\ast(X))$.
For a K3 surface, $H^\ast(X) = H^0(X) \oplus H^2(X) \oplus H^4(X)$ with
$H^0(X) \cong \Q_\ell$ (trivial Galois action) and
$H^4(X) \cong \Q_\ell(-2)$ (Tate twist, $\Frob_p$ acts as $p^2$). The
odd cohomology vanishes because K3 is simply connected with
$h^{1,0} = h^{2,1} = 0$. Summing the three traces yields the stated
identity.
\end{proof}

\noindent\textit{Three computations of Lemma~\ref{lem:phi-k3-bar-shadow-hasse-weil-n1}:}

\begin{description}
\item[Gauss--Jacobi sums.]
Frobenius eigenvalue sums are obtained from explicit Gauss--Jacobi sums
on $\# X(\F_{p^n})$ for $p \in \{3, 5, 7, 11, 13\}$ and
$n \in \{1, 2, 3\}$; see Candelas--de la Ossa--Rodriguez-Villegas
equation (3.12).

\item[Livn\'e and the weight-$3$ CM newform.]
The Fourier coefficients $a_p$ of
$\eta(4\tau)^6 \in S_3(\Gamma_0(16), \chi_{-4})$ give the transcendental
Frobenius trace through
$\Tr(\Frob_p \mid T(X)) = \Tr(\Frob_p \mid H^2) - 20p$.

\item[Hecke Gr\"ossencharacter.]
The singular K3 surface of discriminant $-4$ over $\Q$ is governed by
the Hecke Gr\"ossencharacter $\psi$ of $\Q(i)$ of infinity type $(2,0)$
and conductor $(1+i)^4$.
\end{description}


\section{The proposition}
\label{sec:padic-fermat-proposition}

\begin{proposition}[\texorpdfstring{$p$}{p}-adic CY Langlands for the Fermat quartic]
\label{prop:phi-k3-padic-langlands-fermat}
Let $X \subset \P^3_\Q$ be the Fermat quartic and let $\Phi(X) = H_{\mathrm{Muk}}$ denote its CY-to-chiral image at $d = 2$. For every rational prime $p \neq 2$ of good reduction and every integer $n \geq 1$:

\begin{enumerate}[label=\textup{(\roman*)}]
 \item \textup{(Local-factor identity.)} The $p$-adic shadow zeta $\zeta^{(p)}_{\Phi(X)}$ of $\Phi(X)$, specialised at the Frobenius eigenvalue $\mathbb{L} = p^2$, satisfies the identity of formal Dirichlet series
 \[
 \zeta^{(p)}_{\Phi(X)}(s) \;=\; L_p\bigl(H^2(X, \Q_\ell), s\bigr) \cdot L_p\bigl(H^0(X) \oplus H^4(X), s\bigr),
 \]
 with Dirichlet coefficients
 \[
 a_n(p) \;=\; 1 + \Tr(\Frob_p^n \mid H^2) + p^{2 n}.
 \]
 At $n = 1$ this is Lemma~\ref{lem:phi-k3-bar-shadow-hasse-weil-n1} unconditionally. For $n \geq 2$ the identity requires the full degree-$22$ Frobenius polynomial $P_2(t)$, equivalently the Frobenius action on both $\mathrm{NS}(X)_{\Q_\ell}$ and $T(X)_{\Q_\ell}$. The functional equation constrains this polynomial but does not determine it from the single trace at $n=1$.

 \item \textup{(Livne modularity.)} The $2$-dimensional $\ell$-adic Galois sub-representation on $T(X) \otimes \Q_\ell$ is modular: there exists a weight-$3$ newform
 \[
 f \;=\; \eta(4\tau)^6 \;\in\; S_3\bigl(\Gamma_0(16), \chi_{-4}\bigr)
 \]
 with complex multiplication by $\Q(i)$, such that $a_p(f) = \Tr(\Frob_p \mid T(X))$ for every good prime $p$.

 \item \textup{(Kuga--Satake recovery.)} The polarised Kuga--Satake construction applied to the primitive K3 Hodge structure of rank $21$ gives an abelian variety of dimension $2^{19}$. If one applies the even Clifford construction to the full $H^2$ instead, one obtains the unprimitive variant of dimension $2^{20}$. The Hecke eigensystem of $f$ is carried by the rank-$2$ transcendental factor $T(X)$, equivalently by a CM elliptic factor with multiplication by $\Q(i)$ inside the Kuga--Satake isogeny decomposition; it is not a statement about every factor of the full Kuga--Satake variety.
\end{enumerate}

Clause (i) at $n = 1$ is Lemma~\ref{lem:phi-k3-bar-shadow-hasse-weil-n1}; clauses (ii)--(iii) are ProvedElsewhere via the cited literature.
\end{proposition}

\begin{proof}[Attribution]
Livne (Israel J.~Math.\ 92, 1995) proves the modularity of the $2$-dimensional orthogonal Galois representation on $T(X)$ by a weight-$3$ newform on $\Gamma_0(N)$ for any singular K3 over $\Q$. Sch\"utt (Ramanujan J.\ 19, 2009) identifies the minimal-twist representative for the Fermat quartic as $\eta(4\tau)^6 \in S_3(\Gamma_0(16), \chi_{-4})$, fixing the level $N = 16$ and Nebentypus $\chi_{-4}$. Kuga--Satake (Math.\ Ann.\ 169, 1967) constructs the Clifford abelian variety from the polarised primitive Hodge structure and gives the Galois-equivariant embedding of that primitive part into $\End(H^1(\mathrm{KS}(X)))$; for the Fermat quartic the transcendental rank-$2$ summand is realised by a CM elliptic factor with endomorphism ring $\Z[i]$. Madapusi Pera (Invent.\ Math.\ 201, 2015) underwrites the splitting $H^2 = \mathrm{NS}(1) \oplus T$ as Galois modules. Huybrechts, \emph{Lectures on K3 Surfaces} (CUP 2016), Chapters~4 and~17 consolidate the Hodge--Galois dictionary used here. Clause (i) $n \geq 2$ invokes the full Frobenius characteristic polynomial, constrained by Deligne's Weil bounds but not determined by its trace.
\end{proof}

\paragraph{Scope qualifiers.}
The proposition is specific to the Fermat quartic, not to singular K3 surfaces of Picard rank $20$ in general: other rank-$20$ K3s (for example the $2A_1 + E_8^2$ Inose pencil with discriminant $-3$) have different CM fields and different newform levels. The good-reduction hypothesis excludes $p = 2$, where $X$ has additive reduction and the level-$16$ newform reflects the bad reduction. The standard polarised Kuga--Satake dimension is $2^{19}$ for a K3 surface; the $2^{20}$ number belongs to the full-$H^2$ Clifford variant. The $n \geq 2$ clause of (i) holds only after the degree-$22$ characteristic polynomial of $\Frob_p \mid H^2$ is supplied.


\section{Cross-references and further reading}
\label{sec:padic-fermat-further}

The proposition addresses the $d = 2$ instance of the $p$-adic CY Langlands correspondence. The $d = 3$ analogue for rigid CY threefolds over $\Q$ is the Dieulefait--Manoharmayum theorem (Fields Inst.\ Commun.\ 38, 2003) on Serre modularity; its inclusion is open. The Dwork pencil at $d = 3$ (Candelas--de la Ossa--Rodriguez-Villegas, arXiv:hep-th/0012233) is a candidate for the next family-specific instance. The generic K3 (of Picard rank $< 20$) is \emph{not} covered: the transcendental lattice is then of rank $\geq 3$ and the Galois representation is not $2$-dimensional, so Livne's theorem does not apply directly.

Within Vol~III, Proposition~\ref{prop:phi-k3-padic-langlands-fermat} complements the geometric Langlands correspondence of Chapter~\ref{ch:geometric-langlands}: Feigin--Frenkel gives the chiral centre at the critical level in geometric characteristic, while the present proposition realises the arithmetic analogue at a fixed CM point of K3 moduli. The bridge correspondence $\Phi_2$ is the common input.

\begin{remark}[Arithmetic CY Langlands and the 4d \texorpdfstring{$\mathcal{N}=2$}{N=2} perspective]
\label{rem:padic-4d-n2-arithmetic}
The $p$-adic shadow zeta $\zeta^{(p)}_{\Phi(X)}$ of Proposition~\ref{prop:phi-k3-padic-langlands-fermat}(i) admits a 4d $\mathcal{N}=2$ interpretation at the level of arithmetic. For the Fermat quartic viewed as the Coulomb branch at a CM point, the periods~\eqref{eq:phys-periods} (recorded in the notation of Chapter~\ref{ch:geometric-langlands}, Section~\ref{sec:phys-iib-cy3-hitchin}) specialise to the CM periods $z^I = \oint_{A^I} \Omega$, which are algebraic multiples of the Hecke Gr\"{o}ssencharacter periods of $\Q(i)$. The central charge $Z_\gamma = \oint_\gamma \Omega$ of a D3-brane at the CM point is therefore an algebraic number, not merely a complex period: the BPS states that are massless at the CM point are the arithmetic analogues of the Argyres--Douglas theories of Remark~\ref{rem:phys-ad-irregular}. The Galois action $\rho\colon \mathrm{Gal}(\bar{\Q}/\Q) \to \mathrm{GL}(H^2(X, \Q_\ell))$ of Proposition~\ref{prop:phi-k3-padic-langlands-fermat}(ii) is the arithmetic shadow of the monodromy of the Coulomb branch: in the Hitchin-system language of Section~\ref{sec:phys-iib-cy3-hitchin}, the Frobenius $\mathrm{Frob}_p$ acts on the spectral curve $\Sigma$ by the Hecke correspondence at $p$, and the Livne newform $f = \eta(4\tau)^6$ is the modular form attached to the Hecke eigensheaf on $\Bun_{\mathrm{PGL}_2}(E)$ for the CM elliptic curve $E_i$ (the Kuga--Satake elliptic curve with $j = 1728$). The Seiberg--Witten differential $\lambda_{\mathrm{SW}} = p\,dx$ restricted to $\Sigma$ at the CM point is an algebraic differential, and its periods are the CM periods of~$f$.
\end{remark}
